% Part: normal-modal-logic
% Chapter: completeness
% Section: canonical-models

\documentclass[../../../include/open-logic-section]{subfiles}

\begin{document}

\olfileid[pt]{nml}{com}{cmd}

\olsection{Modelos canônicos}

O \emph{modelo canônico} de um sistema modal~$\Sigma$ é um modelo
específico~$\mModel{M^\Sigma}$ cujos mundos são todos os conjuntos
$\Sigma$-consistentes completos. Sua relação de
acessibilidade~$R^\Sigma$ e sua valoração $V^\Sigma$ são definidas de
modo a garantir que as !!{formula}s verdadeiras em um mundo~$\Delta$
sejam exatamente as !!{formula}s que constituem~$\Delta$.

\begin{defn}
  Seja $\Sigma$ uma lógica modal normal. O \emph{modelo canônico} de
  $\Sigma$ é $\mModel{M}^\Sigma = \tuple{W^\Sigma, R^\Sigma,
    V^\Sigma}$, onde:
  \begin{enumerate}
  \item $W^\Sigma = \Setabs{ \Delta }{\Delta \text{
      é $\Sigma$-consistente completo} }$.
  \item $R^\Sigma \Delta\Delta'$ vale se e somente se
    \iftag{prvBox}{$\Box^{-1}\Delta \subseteq
      \Delta'$}{$\Diamond\Delta' \subseteq \Delta$}.
  \item $V^\Sigma(p) = \Setabs{\Delta}{p \in \Delta}$.
  \end{enumerate}
\end{defn}


\end{document}
